\documentclass[12pt]{article}

\usepackage[a4paper, total={6in, 8in}]{geometry}
\usepackage{textcomp}

\begin{document}
\setlength{\parindent}{0pt}

\def \t {\textrightarrow}
\def \v {\vspace{1em}}
\def \bi {\begin{itemize}}
\def \ei {\end{itemize}}
\def \s[#1] {\section*{#1}}
\def \ss[#1] {\subsection*{#1}}
\def \sss[#1] {\subsubsection*{#1}}

\s[Umberto Saba]
La balia è l'oggetto sostitutivo, ed è il surrogato della madre e del padre \\
Questo sfocia in una nevrosi, che è un aggravazione della psicosi \t si puo associare il tema psicanalitico prima della psicanalisi \\
Lui si sottopone alla terapia del dottor Weiss, che è un allievo di Freud \t segue terapia ma si rende conto che l'inefficacia supera gli sforzi \\
Non trova risoluzione della sua nevrosi \t studiata dai behavioristi, che si differenziano nell'osservazione del soggetto \t un certo comportamento riflette un certo input \\
Dai suoi versi e dalle sue amicizie, c'erano Palazzeschi e altri del gruppo \t forse era omosessuale ?? \\
La poesia onesta \\
Apre il varco alla poetica dell'assenza \t il poeta palombaro con Ungaretti, con anche l'amico del gruppo del cafe di Firenze (dove si trovano tutti gli intellettuali) Eugenio Montale \\
L'ambiente culturale è fervido a trieste \t in tutto questo lui è il cantore del diverso \t ma questo è un valore aggiunto, è il confronto \\
Presenza delle figure familiari, e rapporto edipico con la madre il padre e Peppa \\
Lui ha rispetto solo per la moglie e per la bambina \\

\ss[Città vecchia]
Trieste viene letta attraverso dei luoghi tipici \t non attraverso una perfezione che l'autore non vuole dimostrare \\
Ogni luogo ha un senso, ma non è un luogo edenico \t non è il verdeggiare della vegetazione, non è il sole sulle piazze, ma i luoghi malfamati \t che paradossalmente nascondono la profondità dell animo umano \\
È una sorte di "piazza grande" di Dalla \\
Qua c'è lo sguardo agli ultimi \t verso l'altro, che in tutte le sfaccettature dell'uomo

\v

Spesso = routinarietà \\
Giallo \t contrasto cromatico tra l oscura via e il giallo della pozzanghera \t si ha una relazione da porre con Montale (che è suo amico) \t la il giallo è quello del croco in un polveroso prato \t fiore \\
Il giallo viene dai fanali, che indicano una trieste in fermento \\
La sintassi è molto semplice, aggiunge delle immagini per accumulazione \t no cicerone \\
L'osteria è il luogo di ritrovo, ma il climax ascendente è il lupanare \t è paragonabile alle case di tolleranza \\
Qua ci sono le anti donne angelo \t la stessa angiolina sarebbe stata probabilmente qua \\
Elio vittorini parla di prostituzione, costretta per sfamare i suoi figli \\
Lupanare è termine elettivo, perche viene dalla letteratura latina \\
Uomini e merci sono il detrito \\
Verso 9, io ritrovo l'infinito nell'umiltà \t è qui che c'è la forza sanguigna dell'essere umano \t simile a Pier Paolo Pasolini, la cultura borghese vede bontà nell'etichetta con pregiudizi \\
Federico Borromeo, che abbraccia i ragazzi poveri \\
La cultura borghese è limitante \t qua Saba da un altro volto \\
Presenta figure di un umanità palpitante \t dragone, nella tradizione mercenaria del 600 e del 700 \t espressioni nuove calate in un'umanità palpitante, senza filtri \\
Tumultuante, impeto dell'anima, che perde il controllo e fa emerge la visceralita \\
Creatura della vita, e a capo "del dolore" \t isolato del dolore \t e parlare di vita significa parlare di dolore \\
Non è necessariamente cristiano \\
Il suo pensiero si fa piu puro, dove la strada presenta immagini turpi e disdicevoli \\
Ultimo verso, puro e tupre creano un incrocio (quasi un chiasmo) \\
Fonti editoriali sono mondadori e einaudi \\
Emblema della città, e lirica \t citta fumose con case aggiunte a case \t la città letta nello Spleen \t o anche Coketown, e l'opposto della città dei cesari del piacere \t la città di dalla la piazza grande, anche qua umanita \\
Riqualificazione di alcune aree di periferia degradate \\
In verga è diverso 

\v

De Andre parla anche di questo e compone canzone omonima \\
Terenzio si ricollega anche

\ss[Mio padre è stato per me un assassino]
La tenzone è l'artifizio letterario per cui ?? \\
La madre è Rachele Coenne (di origine ebrea) \t il padre è pellegrino perche ha abitudine a mouversi di continuo \\
L'assassino tra virgolette, e l'articolo evidenzia che è il vero assassino \\
Descrive quando l'ha conosciuto e per quanto, e come fosse in realtà ancora immaturo \\
Il dono + grande che ha ricevuto è lo sguardo \t sguardo azzurrino = metonimia, il sorriso dolce e astuto (sinestesia e non veramente ossimoro, dolce puo essere ottenere altro) \\
Questo suo cercare si traduce in passare da donna in donna \t pasciuto da pasco, che vuol dire nutrito e alimentato \\
Leggerezza calviniana è quella di trattare temi impegnativi guardandoli da un altra prospettiva \\
Il pallone è qualcosa che è caratterizzato dalla giocosità sportiva, ma anche quella dei bambini \t come se il padre non fosse mai crescuto \\
Gli è sfuggito di mano \t una pesantezza perche era da educare, leggerezza perchè comunque era come un bambino quindi non gli si possono dare colpe vere \\
Linguaggio colloqiuale, senso dell'abbandono, figura di una donna descritta come assertiva (lei cerca di trattenerlo ma gli sfugge ed è rigorosa con il figlio, ammonisce con imperativo) \\
Quando l'autore si dedica alla autonalisi e smette di seguire terapia \\
Si cogli eimmediatezza a parlare anche di se \t la sua tempra è piu votata a parlare di se che degli altri, a scandaglare il suo animo \t anticipa psicanalisi \\
Mostra tempra empatica ed emotiva

\ss[Dico al mio cuore]
Da vedere in parallelo "A se stesso" di Leopardi \t è una confessione in prima persona, e mostra il suo mondo interiore \\
La moglie lo abbandona per un periodo \\
Celebrazione del non cambiamento \t lui non puo non amarla \\
Parla con il suo cuore, quindi autoanalisi \t scordala, voce del cuore e voce della ragione \\
Ti vedo \t tutto si cancella \\
Perchè generoso e vile? \t generoso perchè da amore a lei, ma non ha il coraggio di rifiutarla \\
Non c'è la stessa dimensione del cuore lacerato, come in Leopardi \\
Lei non coglie neanche il dolore che causa in lui per l'abbandono \\
La sua è una voce amara = sinestesia \t non c'è tessuto retorico fitto, entra in ottica l'epopea del quotidiano
\end{document}
